% Source: https://artofproblemsolving.com/community/c3h2567808_nasty_trig_prob
% Stable identity: aops-topic-2567808-q01
Write $s=\sin\theta$ and $c=\cos\theta$. Rotating the point $(c,s)$ of the unit circle by $90^\circ$ gives $(-s,c)$, by $180^\circ$ gives $(-c,-s)$, by $270^\circ$ gives $(s,-c)$, and reflecting it in the $x$-axis gives the point at angle $-\theta$ (or $360^\circ-\theta$), namely $(c,-s)$. Hence
\begin{align*}
&\sin(180^\circ+\theta)=-s,\quad \cos(90^\circ+\theta)=-s,\quad \sin(360^\circ-\theta)=-s,\quad \cos(360^\circ+\theta)=c,\quad \csc(-\theta)=-\dfrac1s,\\
&\tan(270^\circ-\theta)=\dfrac{\sin(270^\circ-\theta)}{\cos(270^\circ-\theta)}=\dfrac{-c}{-s}=\dfrac cs,\qquad \cot(360^\circ-\theta)=\dfrac{c}{-s}=-\dfrac cs,\\
&\cot(270^\circ+\theta)=\dfrac{\cos(270^\circ+\theta)}{\sin(270^\circ+\theta)}=\dfrac{s}{-c}=-\dfrac sc,\qquad \sin(270^\circ+\theta)=-c.
\end{align*}
(For $270^\circ-\theta$: the point at angle $-\theta$ is $(c,-s)$, and rotating it by $270^\circ$ gives $(-s,-c)$.)

\textbf{(a)} The numerator is $(-s)(-s)\cdot\dfrac cs\cdot\left(-\dfrac cs\right)=-c^2$, and the denominator is
\[(-s)\cdot c\cdot\left(-\dfrac1s\right)\cdot\left(-\dfrac sc\right)=-s.\]
So the left-hand side equals $\boxed{\dfrac{\cos^2\theta}{\sin\theta}}$, which is not identically $1$: for $\theta=30^\circ$ it equals $\dfrac{3}{4}\div\dfrac12=\dfrac32$. The claim is false.

\textbf{(b)} With $\sin(270^\circ+\theta)=-c$ in place of $\cot(270^\circ+\theta)$, the denominator becomes $(-s)\cdot c\cdot\left(-\dfrac1s\right)\cdot(-c)=-c^2$, so the left-hand side is $\dfrac{-c^2}{-c^2}=1$ whenever it is defined. $\blacksquare$

\begin{asy}
usepackage("amsmath");
size(5cm);
defaultpen(fontsize(10pt));
real t=35;
draw(circle((0,0),1));draw((-1.3,0)--(1.3,0),Arrow(TeXHead));draw((0,-1.3)--(0,1.3),Arrow(TeXHead));
pair P=dir(t),Q=dir(t+90),R=dir(t+180),T=dir(-t);
draw((0,0)--P);draw((0,0)--Q,dashed);draw((0,0)--R,dashed);draw((0,0)--T,dashed);
dot(P);dot(Q);dot(R);dot(T);
draw(arc((0,0),0.25,0,t));label("$\theta$",0.37*dir(t/2),fontsize(8pt));
label("$\theta$",P,NE,fontsize(9pt));label("$90^\circ+\theta$",Q,NW,fontsize(9pt));
label("$180^\circ+\theta$",R,SW,fontsize(9pt));label("$360^\circ-\theta$",T,SE,fontsize(9pt));
\end{asy}

The sketch shows the points of the unit circle at angles $\theta$, $90^\circ+\theta$, $180^\circ+\theta$ and $360^\circ-\theta$ used above.
